\documentclass[twocolumn]{article}

%%%%%%%%%%%%%%%%%%%%%%%%%%%%%%%%%%%%%%%%%%%%%%%%%%%%%%%%%%%%%%%%%%
% Packages
\usepackage[T1]{fontenc}
\usepackage[utf8]{inputenc}
% --- FONTS ---
\usepackage{tgtermes} % Main text font: Times New Roman style
\usepackage{tgheros}  % Sans-serif font: Helvetica/Arial style (for bold headers)
\usepackage{microtype} 

\usepackage{graphicx}
\usepackage{tabularx}
\usepackage{ragged2e}
\usepackage{booktabs}
\usepackage{amsmath}
\usepackage{amsfonts}
\usepackage{amssymb}
\usepackage{longtable}
\usepackage[table]{xcolor}
\usepackage{caption}
\usepackage{hyperref}
\usepackage{csquotes}
\usepackage{geometry}
\usepackage{fancyhdr}
\usepackage{titlesec}
\usepackage{float}
\usepackage{lipsum} 
\usepackage{listings}

% --- Bibliography (Strictly APA 7th) ---
\usepackage[style=apa, backend=biber]{biblatex}
\addbibresource{references.bib} % Connects to your references.bib file

% --- Custom Styling ---
\geometry{lmargin=0.6in,rmargin=0.6in,tmargin=0.75in,bmargin=0.75in,footskip=20pt,columnsep=0.3in}

% Headers
\pagestyle{fancy}
\fancyhf{}
\fancyhead[L]{<Course Name>}
\fancyhead[R]{\today}
\fancyfoot[C]{\thepage}

\graphicspath{ {images/} }

% Section Titles: Robust calculation to prevent margin bleed
\newcommand{\styledtitle}[1]{%
  \noindent\colorbox{black}{\parbox{\dimexpr\linewidth-2\fboxsep\relax}{\centering\textcolor{white}{\sffamily\bfseries\MakeUppercase{#1}}}}%
}

\titleformat{\section}
  {\normalfont}{}{0em}{\styledtitle}

\titleformat{\subsection}
  {\normalfont\normalsize\sffamily\bfseries}{}{0em}{}

\titleformat{\subsubsection}
  {\normalfont\normalsize\sffamily\itshape}{}{0em}{}

\lstset{
    language=Python, 
    basicstyle=\ttfamily\scriptsize, 
    breaklines=true,
    breakatwhitespace=true,
    frame=single,
    backgroundcolor=\color{gray!5},
    keywordstyle=\color{blue}\bfseries,
    commentstyle=\color{green!40!black},
    stringstyle=\color{red},
    numbers=left,
    numberstyle=\tiny\color{gray},
    stepnumber=1,
    showstringspaces=false,
    showlines=true
}

\hypersetup{
    colorlinks=true,
    linkcolor=black,
    filecolor=black,      
    urlcolor=blue,
    citecolor=black
}

%%%%%%%%%%%%%%%%%%%%%%%%%%%%%%%%%%%%%%%%%%%%%%%%%%%%%%%%%%%%%%%%%%

\begin{document}

%%%%%%%%%%%%%%%%%%%%%%%%%%%%%%%%%%%%%%%%%%%%%%%%%%%%%%%%%%%%%%%%%%
%%%%%%%%%%%%%%% AUTHORS AND TITLE INFO %%%%%%%%%%%%%%%
%%%%%%%%%%%%%%%%%%%%%%%%%%%%%%%%%%%%%%%%%%%%%%%%%%%%%%%%%%%%%%%%%%
\twocolumn[
\begin{@twocolumnfalse}
    \begin{center}
        \LARGE \textbf{<Template Title: Your Academic Article or Diploma Design>} \\ \vspace{0.5cm}
        
        \normalsize \textbf{Amangeldy Shyngyskhan\textsuperscript{1}, <Second Author>\textsuperscript{1}} \\ \vspace{0.1cm}
        \normalsize \textit{\textsuperscript{1}Perricheno Team | Astana IT University} \\ \vspace{0.4cm}
        
        % Flattened horizontally to save space and unified font sizes:
        \normalsize \textbf{Course:} <Course Name> \qquad \textbf{Group:} <Your Group> \qquad \textbf{Supervisor:} <Name> \\ \vspace{0.2cm}
        
        \normalsize \today \ | Project Report / Diploma Design
    \end{center}
    \vspace{0.3cm}
    
    \noindent\textbf{Abstract} \\
    \textit{This study provides a comprehensive quantitative analysis of the chosen dataset. Please replace this abstract text with a concise summary of your research, covering the background, methodology, key findings, and conclusion. Ensure it is approximately 150-250 words. Note that the font here is safely italicized across all text.}
    \vspace{0.6cm}
\end{@twocolumnfalse}
]

%%%%%%%%%%%%%%%%%%%%%%%%%%%%%%%%%%%%%%%%%%%%%%%%%%%%%%%%%%%%%%%%%%
%%%%%%%% MAIN TEXT %%%%%%%%%
%%%%%%%%%%%%%%%%%%%%%%%%%%%%%%%%%%%%%%%%%%%%%%%%%%%%%%%%%%%%%%%%%%

\section{INTRODUCTION}
\lipsum[2]

\section{LITERATURE REVIEW}
\lipsum[3]

% Example of citing placeholder (Using APA 7th and entries from your references.bib)
\noindent According to \textcite{nguyen2018impact}, delivery speed significantly impacts customer satisfaction. Another study confirms this trend \parencite{rao2011}.

\section{DATA EXPLORATION}
\lipsum[4]

\begin{itemize}
    \item \textbf{Key Point A}: Analysis of major dataset factors.
    \item \textbf{Key Point B}: Distribution of logistics components.
\end{itemize}

\section{METHODOLOGY}
\lipsum[5]

\subsection{Statistical Summary}
We present the correlation structure of our variables in Table \ref{tab:template_table}. 

\begin{table}[H]
\centering
\caption{Example Table Template}
\label{tab:template_table}
\resizebox{\columnwidth}{!}{%
\begin{tabular}{lcccc}
\toprule
 & Column A & Column B & Column C & Column D \\
\midrule
Row 1 & 1.00 & -0.28 & 0.08 & 0.07 \\
Row 2 & -0.28 & 1.00 & -0.09 & 0.24 \\
Row 3 & 0.08 & -0.09 & 1.00 & 0.52 \\
\bottomrule
\end{tabular}%
}
\end{table}

\section{RESULTS \& DISCUSSION}

\subsection{1. First Key Finding}
\textbf{Business Target}: Placeholder analysis output here.

\textbf{Extended Analysis}: Placeholder explanation output here.

\begin{lstlisting}[caption={Example Code Snippet}]
# Placeholder for data analysis code
import pandas as pd
import numpy as np

def run_analysis(data):
    """
    Function to compute core metrics
    """
    return np.mean(data)
\end{lstlisting}

\begin{figure}[H]
    \centering
    \fbox{\rule{0pt}{1.5in} \rule{0.9\linewidth}{0pt}} 
    \caption{Placeholder for an Example Data Figure}
    \label{fig:example_plot}
\end{figure}

\subsection{2. Second Key Finding}
\textbf{Impact}: Key impacts detailed here.

\textbf{Extended Analysis}: Details and correlations.

\begin{figure}[H]
    \centering
    \fbox{\rule{0pt}{1.5in} \rule{0.9\linewidth}{0pt}} 
    \caption{Placeholder for Another Example Figure}
    \label{fig:example_plot2}
\end{figure}

\section{CONCLUSION}
\lipsum[10]

% -------------------------------------------------------------
% BIBLIOGRAPHY
% -------------------------------------------------------------
\nocite{*} % Prints ALL references from your references.bib file, even if you didn't cite them!
\printbibliography

\end{document}
